% !TEX program = lualatex
% =====================================================================
% execucao-desempenho.tex -- nota tecnica
%
% Razao de desempenho da frota fotovoltaica, separando o que e decisao de
% despacho do que e do ativo.
%
% Nao entra no `reporte.tex`. Compila sozinho:
%   latexmk -pdf execucao-desempenho.tex
% =====================================================================

% A folha de estilo pressupoe `book`: usa `\chaptertitlename` e os
% comandos de capitulo do `tocloft`. Em `article` ela nao carrega.
\documentclass[10pt,a4paper,oneside,twocolumn]{book}

\usepackage{iftex}
\ifPDFTeX\usepackage[utf8]{inputenc}\fi
\usepackage[brazilian]{babel}
\usepackage[a4paper,top=2.0cm,bottom=2.2cm,left=1.9cm,right=1.9cm,
            columnsep=0.7cm]{geometry}
\usepackage{estilo-reporte}

\raggedbottom
\pagestyle{plain}
\setlength{\parskip}{0.35em}
\setlength{\parindent}{0pt}

\marca{../docs/img/terra-mark.png}
\autoria{João Leonardi da Silva Melo}{joao\_leonardi.melo@somosicev.com}

\begin{document}
\abertura{Razão de desempenho da frota fotovoltaica}
         {Nota técnica · \today{} ·
          \texttt{notebooks/br/\_desempenho.py}}

\section*{Objetivo}

Decompor o déficit de desempenho da frota fotovoltaica em decisão de despacho e
desempenho do ativo.

A razão de desempenho é energia entregue sobre energia que a irradiância medida
no plano permitiria, $\mathrm{PR}=(E/P_{\mathrm{nom}})/(\mathrm{POA}/1000)$. O
denominador depende do plano do arranjo, e por isso o cálculo requer a
classificação de geometria da nota anterior.

\section*{Dados}

Janela de 2025. Geração e irradiância do \texttt{br.pv\_detail}, por usina;
corte e disponibilidade do \texttt{br.pv\_curtail}, por conjunto. A ligação
entre os dois é por nome normalizado, e casa 96\% dos conjuntos.

\begin{ficha}
  \item[Amostra] \num{69297} dias-usina em 236 usinas com geometria
                 classificada, das quais 235 têm 60 dias ou mais.
\end{ficha}

\campo{Não há indisponibilidade declarada de fotovoltaica} O
\texttt{taxa\_teif\_teip} cobre hidrelétrica, térmica e nuclear. O
\texttt{disponibilidade\_usina} devolve zero linhas de UFV. A única
disponibilidade de fotovoltaica publicada está dentro do próprio
constrained-off, por conjunto e não por usina.

\section*{Método}

Três quantidades por usina, medianas dos seus dias:

\begin{ficha}
  \item[Realizado] com a geração verificada.
  \item[Potencial] com a geração estimada pelo ONS, que é o que a usina teria
                   entregue sem restrição segundo o modelo do operador.
  \item[Sem corte] realizado nos dias sem intervalo restrito no conjunto.
\end{ficha}

\campo{Controle} Somando as usinas de um conjunto, a geração do
\texttt{pv\_detail} tem de igualar a do \texttt{pv\_curtail}. Medido em
\num{6102} conjunto-dias com todas as usinas presentes: razão 1,0000.

\section*{Resultados}

\vspace{0.2em}
\tabfont
\begin{tabular}{@{}>{\rag\arraybackslash}p{3.2cm}rrr@{}}
\toprule
\cab{Razão de desempenho} & \cab{p10} & \cab{Mediana} & \cab{p90} \\
\midrule
Realizado        & 0,551 & \textbf{0,838} & 0,961 \\
Potencial (ONS)  & 0,864 & 1,023 & 1,208 \\
Sem dia restrito & 0,858 & 0,998 & 1,131 \\
\bottomrule
\end{tabular}
\notas{Mediana entre 235 usinas, cada uma mediana dos seus dias de 2025.}

\vspace{0.6em}
A lacuna entre o potencial e o realizado é 0,227. Restringindo aos dias sem
restrição no conjunto, ela cai para 0,031. Cerca de 86\% do déficit é decisão de
despacho. Dias com alguma restrição no conjunto são 80,7\% do total.

\campo{Contabilidade de energia} Nos \num{5626} conjunto-dias com todas as
usinas presentes, 46 conjuntos:

\vspace{0.3em}
\tabfont
\begin{tabular}{@{}>{\rag\arraybackslash}p{4.4cm}r@{}}
\toprule
\cab{Energia, 2025} & \cab{GWh} \\
\midrule
Estimada pelo ONS       & \num{11610.9} \\
Verificada              &  \num{8016.0} \\
Lacuna                  &  \num{3594.8} \\
Corte medido            &  \num{3610.7} \\
\bottomrule
\end{tabular}

\vspace{0.5em}
O corte medido responde por 100\% da lacuna. Estimativa e medida de corte se
sustentam uma contra a outra.

\campo{Por geometria}
\vspace{0.2em}

\tabfont
\begin{tabular}{@{}>{\rag\arraybackslash}p{2.5cm}rrr@{}}
\toprule
\cab{Geometria} & \cab{$n$} & \cab{Realiz.} & \cab{Potenc.} \\
\midrule
Fixo        &  93 & 0,901 & 1,154 \\
Rastreador  & 142 & 0,795 & 0,971 \\
\bottomrule
\end{tabular}

\section*{Limitações}

\campo{A separação de três vias reduz-se a duas} Sem indisponibilidade
declarada, o resíduo de 0,031 mistura subdesempenho do ativo com
indisponibilidade não declarada. É um teto para o subdesempenho, não uma medida
dele.

\campo{Não é razão de desempenho de norma} Está referenciada à potência
fiscalizada, em corrente alternada; a potência de corrente contínua dos módulos
não é publicada pela ANEEL. Um arranjo é sobredimensionado em CC contra o
inversor, e por isso o índice passa de~1.

\campo{O modelo do ONS não distingue a montagem} O PR potencial mediano é 1,154
nas fixas contra 0,971 nas rastreadoras, resultado esperado de um modelo que
supõe arranjo fixo. Razão CC/CA menor em usina com rastreador prevê o mesmo
sinal, e a potência CC não é publicada.

\begin{objecao}
As duas explicações do parágrafo acima não se separam com o dado disponível. O
sinal fica registrado como sinal.
\end{objecao}

\section*{Perfis fora da distribuição}

O platô diurno mediano da frota é 0,903. Três usinas ficam abaixo de 0,75, e as
duas mais afastadas têm defeito identificado no perfil de meia em meia hora:
Guaimbé~1 com depressão no meio do dia entre dois máximos, compatível com
sombreamento parcial sobre o sensor, e Santa Luzia~4 com a leitura zerando às
13h00 e não retornando.

\section*{Artefatos}

\begin{ficha}
  \item[Código] \texttt{notebooks/br/\_desempenho.py}.
  \item[Saída]  \texttt{data/processed/desempenho\_por\_usina.csv}.
\end{ficha}

\end{document}
